% Source: https://artofproblemsolving.com/community/c3h1849294_geometry
% Stable identity: aops-topic-1849294-q01
\textbf{Construction.} Construct the midpoint $M$ of $O_1O_2$, draw the segment $AM$, and construct the line $\ell$ through $A$ perpendicular to $AM$. Let $\ell$ meet the first circle again at $B$ and the second at $C$. Then $AB=AC$.

\begin{asy}
usepackage("amsmath");
size(7cm);
defaultpen(fontsize(10pt));
pair O1=(0,0),O2=(4.2,0.8);real r1=3,r2=2.3;
pair[] I=intersectionpoints(circle(O1,r1),circle(O2,r2));
pair A=(I[0].y>I[1].y)?I[0]:I[1];
pair M=(O1+O2)/2;
pair u=unit(rotate(90)*(A-M));
pair B=A+2*dot(O1-A,u)*u,C=A+2*dot(O2-A,u)*u;
pair F1=A+dot(O1-A,u)*u,F2=A+dot(O2-A,u)*u;
draw(circle(O1,r1),linewidth(0.9));draw(circle(O2,r2),linewidth(0.9));
draw((B+0.5*(B-C)/abs(B-C))--(C+0.5*(C-B)/abs(C-B)),linewidth(0.6));
draw(B--C,red+linewidth(1.3));
draw(O1--O2,dashed);draw(A--M,blue+dashed);
draw(O1--F1,dotted);draw(O2--F2,dotted);
void ra(pair V,pair d1,pair d2){real s=0.25;draw(V+s*d1--V+s*d1+s*d2--V+s*d2,linewidth(0.5));}
ra(A,unit(M-A),u);ra(F1,unit(O1-F1),u);ra(F2,unit(O2-F2),-u);
pair n=unit(rotate(90)*u);
pair[] mids={(A+F1)/2,(F1+B)/2,(A+F2)/2,(F2+C)/2};
for(pair q:mids) draw(q-0.15*n--q+0.15*n,red+linewidth(0.8));
dot(A);dot(B);dot(C);dot(O1);dot(O2);dot(M,blue);dot(F1);dot(F2);
label("$A$",A,2.5*N);label("$B$",B,dir(110));label("$C$",C,dir(60));
label("$O_1$",O1,SW);label("$O_2$",O2,SE);label("$M$",M,dir(-30),blue);
label("$F_1$",F1,N);label("$F_2$",F2,N);
label("\mbox{}",(min(currentpicture,user=true)+max(currentpicture,user=true))/2,black); // keeps the text after the figure black when LaTeX typesets the labels inline
\end{asy}

\textbf{Proof.} Let $F_1$ and $F_2$ be the feet of the perpendiculars from $O_1$ and $O_2$ to $\ell$. The perpendicular from the centre of a circle to a chord bisects the chord, so $F_1$ is the midpoint of $AB$ and $F_2$ is the midpoint of $AC$; hence $AB=2AF_1$ and $AC=2AF_2$.

The lines $O_1F_1$, $MA$ and $O_2F_2$ are all perpendicular to $\ell$, so they are parallel. In the trapezoid $O_1F_1F_2O_2$ (or along any line, after projecting onto $\ell$), the point $M$ is the midpoint of $O_1O_2$, so its projection $A$ onto $\ell$ is the midpoint of $F_1F_2$, by the midline theorem (parallel lines cutting equal segments on one transversal cut equal segments on another). Therefore $AF_1=AF_2$, and so
\[
AB=2AF_1=2AF_2=AC .
\]
The points $F_1$ and $F_2$ are on opposite sides of $A$, so the chords $AB$ and $AC$ lie on opposite sides of $A$, as required. (The line $\ell$ is not tangent to either circle: it would be tangent to the first circle only if $\ell\perp AO_1$, i.e.\ if $M$ lay on line $AO_1$, which would put $O_2$ on that line too and make the circles tangent at $A$.)

\textbf{Uniqueness.} Conversely, if a line through $A$ cuts equal chords $AB=AC$ on opposite sides of $A$, then $A$ is the midpoint of $F_1F_2$, so it is the projection of $M$ onto the line, i.e.\ the line is perpendicular to $AM$. So $\ell$ is the only such line. (If the chords may lie on the same side of $A$, the common chord of the two circles is a second, trivial solution: there the two chords coincide.)
